% =============================================================
% selfwork_landscape.tex — Шаблон самостоятельной работы
% Формат: A4 landscape, 2 варианта на одном листе
% Компилятор: XeLaTeX (два прохода для линий реза)
%
% Раскладка:
%   ┌───────────────────┊───────────────────┐
%   │    Вариант 1      ┊    Вариант 2       │
%   │                   ┊                   │
%   │  [шапка]          ┊  [шапка]          │
%   │  Имя и фам.:___   ┊  Имя и фам.:___   │
%   │                   ┊                   │
%   │  1. ...           ┊  1. ...           │
%   │  2. ...           ┊  2. ...           │
%   │  ...              ┊  ...              │
%   │                   ┊                   │
%   └───────────────────┊───────────────────┘
%                       ↑ вертикальная линия реза
% =============================================================

\documentclass[12pt]{article}
\usepackage{styles/mathworksheet}

% ============================================================
% ПАРАМЕТРЫ РАБОТЫ — ИЗМЕНЯЙТЕ ТОЛЬКО ЗДЕСЬ
% ============================================================
\newcommand{\worktitle}{Самостоятельная работа}
\newcommand{\worktopic}{Линейные уравнения}
% ============================================================

\begin{document}

% ============================================================
% ЗАДАНИЯ — определяются как макросы отдельно от структуры
% ============================================================

% ---- ВАРИАНТ 1 ----------------------------------------------
\newcommand{\variantA}{%
  % --- Блок 1: решить уравнения (короткий формат) ---
  \calcrowc{1}{3x - 7 = 14}
  \calcrowc{2}{2(x + 5) = 16}
  \calcrowc{3}{\tfrac{x}{4} - 3 = 1}
  \tasksep
  % --- Блок 2: задание со свободным форматом ---
  \taskrow{4}{Решите уравнение и проверьте ответ:\quad
    $5x + 3 = 2x - 9$}
  \vspace{12mm}%
  \tasksep
  % --- Блок 3: задание-выражение ---
  \taskrow{5}{При каком значении $x$ выражение
    $\dfrac{x - 2}{3}$ равно нулю?}
  \vspace{8mm}%
  \tasksep
  % --- Блок 4: задача ---
  \taskproblem{6}{%
    Поезд прошёл первую часть пути за $3$ ч со скоростью
    $80$ км/ч, а вторую~— за $2$ ч со скоростью $x$ км/ч.
    Весь путь равен $400$ км. Найдите $x$.%
  }{25}
}

% ---- ВАРИАНТ 2 ----------------------------------------------
\newcommand{\variantB}{%
  \calcrowc{1}{4x + 9 = 21}
  \calcrowc{2}{3(x - 4) = 15}
  \calcrowc{3}{\tfrac{x}{3} + 2 = 6}
  \tasksep
  \taskrow{4}{Решите уравнение и проверьте ответ:\quad
    $7x - 5 = 4x + 10$}
  \vspace{12mm}%
  \tasksep
  \taskrow{5}{При каком значении $x$ выражение
    $\dfrac{x + 5}{4}$ равно нулю?}
  \vspace{8mm}%
  \tasksep
  \taskproblem{6}{%
    Велосипедист проехал первую часть пути за $2$ ч со
    скоростью $15$ км/ч, а вторую~— за $x$ ч со скоростью
    $12$ км/ч. Весь путь равен $54$ км. Найдите $x$.%
  }{25}
}

% ============================================================
% ДВУХКОЛОНОЧНЫЙ РАЗВОРОТ
% ============================================================
\begin{variantpage}
  \swblock{\worktitle}{\worktopic}{1}{\variantA}
  &
  \swblock{\worktitle}{\worktopic}{2}{\variantB}
\end{variantpage}

\end{document}

% ============================================================
% СПРАВОЧНИК КОМАНД
% ============================================================
%
%  Шапка и структура:
%    \swblock{название}{тема}{вариант}{задания}
%      — полный блок варианта (используйте внутри variantpage)
%
%  Задания:
%    \calcrowc{N}{выр}      — вычислит. задание (textstyle)
%    \calcrow{N}{выр}       — вычислит. задание (displaystyle)
%    \taskrow{N}{текст}     — произвольное задание
%    \taskrowans{N}{текст}  — задание + строка ответа
%    \taskproblem{N}{условие}{высота_мм}  — задача с полем
%    \tasksep               — лёгкий разделитель
%    \ansline[ширина]       — линия для ответа
%    \mf{формула}           — строчная формула (textstyle)
%
%  Для коротких заданий (7–8 штук):
%    Используйте \calcrowc — плотный формат
%    Добавьте \tasksep между смысловыми блоками
%
%  Для смешанного контента (4–6 заданий с пространством):
%    Используйте \taskrow + \vspace{...mm} для воздуха
%    Финальная задача через \taskproblem{}{}{высота}
%
%  КОМПИЛЯЦИЯ:
%    xelatex selfwork_landscape.tex
%    xelatex selfwork_landscape.tex   ← второй проход обязателен
% ============================================================




